% ============================================================
%  ch04c：§4.5 定理 3.5.2.1 证明的设定、§4.6 L-值与迹
%  （原文 pp.31 中段–39，起于 §4.5，止于引理 4.6.6.4 的证明）
% ============================================================
\subsection{定理 3.5.2.1 证明的设定}

我们开始证明定理 3.5.2.1。

\pnum{回忆我们的目标是定义交换图 (3.11) 并验证它的若干性质。}

我们首先定义函数 $\delta\colon\bT\to\A^1$ 如下：
\[
  \delta(t) := \prod_{\lambda\in\mathcal{R}_X}\ch_{\bg}(t,\lambda),\qquad t\in\bT,
\]
其中我们使用 §4.3.2 的记号，并把 $\bg$ 通过伴随作用看作 $\bT$ 的表示。更显式地，我们有：
\[
  \delta(t) = \prod_{\lambda\in\mathcal{R}_X}(1-\lambda)^{\dim(\bT)}\prod_{\check\alpha\in\check\Delta}(\check\alpha(t)-\lambda),
\]
其中我们把 $\check\alpha$ 看作映射 $\bT\to\mathbb{G}_m\subset\A^1$。

显然 $\delta$ 是 $W$-不变态射，因此诱导一个映射 $\bT/\!/W\to\A^1$；我们仍用 $\delta$ 记这个函数。

\pnum{现在我们构造一个交换图：}
\begin{equation}
\begin{tikzcd}[column sep=0.9em, row sep=1.5em]
& \LS^{\arthm,\circ}_{\bB} \arrow[dl] \arrow[d, "\chi_{\bB}"] \arrow[dr] & \\
\LS^{\arthm,\circ}_{\bT} \arrow[d, "\chi_{\bT}"]
  & \bB^{\ad}/\bB \arrow[dl] \arrow[dr]
  & \LS^{\arthm,\circ}_{\bG} \arrow[d, "\chi_{\bG}"] \\
\bT^{\ad}/\bT \arrow[dr] & & \bG^{\ad}/\bG \arrow[dl] \\
& \bT/\!/W \arrow[d, "\delta"] & \\
& \A^1 &
\end{tikzcd}
\end{equation}

我们现在这样定义 $\delta_{\bG}$：它在 $\LS^{\arthm,\circ}_{\bG}$ 上的限制由 (4.8) 中出现的唯一映射 $\LS^{\arthm,\circ}_{\bG}\to\A^1$ 给出，而它在 $\LS^{\arthm}_{\bG}\setminus\LS^{\arthm,\circ}_{\bG}$ 上的限制\footnote{我们提醒，
$\LS^{\restr,\circ}_{\bG}\subseteq\LS^{\restr}_{\bG}$ 是一个连通分支，因此先在
$\LS^{\arthm,\circ}_{\bG}:=\LS^{\arthm}_{\bG}\times_{\LS^{\restr}_{\bG}}\LS^{\restr,\circ}_{\bG}$
上定义这个函数、并在别处令其为零，是合法的。}恒为零。我们用完全同样的方式定义 $\delta_{\bT}$，只需把上一句中的 $\bG$ 处处换成 $\bT$。

下面我们验证映射对 $(\delta_{\bG},\delta_{\bT})$ 满足定理 3.5.2.1 的结论。

\pnum{首先，由 (4.8) 显然可知交换图 (3.11) 存在。}

\pnum{其次，我们需要验证定理 3.5.2.1 的 (1)，即 $\delta_{\bG}$ 在平凡 Weil 局部系统处取非零值。由构造，只需证明 $\delta(1)\neq 0$。显然
$\delta(1)=\prod_{\lambda\in\mathcal{R}_X}(1-\lambda)^{\dim(G)}$，而我们记得 $1\notin\mathcal{R}_X$（见注 4.3.1.1）。}

\pnum{接下来我们注意到，$\delta_{\bG}$ 非零的轨迹恰好是 $\LS^{\arthm,\circ,\nonres}_{\bG}$（按定义）。因此，由定理 4.3.3.1，我们有：}
\begin{equation}
  \Gamma(\LS^{\arthm}_{\bG},\omega)[\delta_{\bG}^{-1}]
  = \Gamma(\LS^{\arthm,\circ,\nonres}_{\bG},\omega)
  \cong \Gamma((\bG^{\ad}/\bG)^{\nonres},\omega),
\end{equation}
它在次数 $0$ 集中，因为它是 $\Gamma(\bG^{\ad}/\bG,\omega)\cong\Gamma(\bG^{\ad}/\bG,\OO)$ 在函数
$\bG^{\ad}/\bG\to\A^1$ 处的局部化，从而验证了定理 3.5.2.1 的假设 (2)。

\pnum{剩下要验证满射性（即定理 3.5.2.1 的 (3)）。我们将在本节余下部分做这件事；这里先作一些预备性的、起定向作用的评注。}

回忆引理 3.7.1.1 的设定。我们注意到有两个映射：
\[
\begin{gathered}
\Gamma(\LS^{\arthm}_{\bT},\omega)[\delta_{\bT}^{-1}]
\xrightarrow{\Eisspec}
\Gamma(\LS^{\arthm}_{\bG},\omega)[\delta_{\bG}^{-1}]
= \Gamma(\LS^{\arthm,\circ,\nonres}_{\bG},\omega)\\
\stackrel{(4.9)}{\cong}
\Gamma(\bG^{\ad}/\bG,\omega)[(\delta_{\bG}\circ\tau_{\bG})^{-1}]
= \Gamma((\bG^{\ad}/\bG)^{\nonres},\omega),
\end{gathered}
\]
以及：
\[
\begin{gathered}
\Gamma(\LS^{\arthm}_{\bT},\omega)[\delta_{\bT}^{-1}]
\stackrel{\text{定理 4.3.3.1}}{=}
\Gamma(\bT^{\ad}/\bT,\omega)[(\delta_{\bT}\circ\tau_{\bT})^{-1}]
\xrightarrow{\mathrm{Eis}^{\mathrm{spec},\mathrm{toy}}}
\Gamma(\bG^{\ad}/\bG,\omega)[(\delta_{\bG}\circ\tau_{\bG})^{-1}]\\
= \Gamma((\bG^{\ad}/\bG)^{\nonres},\omega).
\end{gathered}
\]
由推论 3.7.2.1，如果这两个映射重合，我们就完成了。

这一期望有点太天真：我们转而证明它们相差可逆的 $L$-值，这对我们的目的来说已经足够。

\subsection{$L$-值与迹}

在 §4.5.6 中，我们对 $L$-值作了一个晦涩的评注。在本小节中，我们将在范畴迹与 $L$-值之间建立精确的联系；这是我们所需的主要计算输入。

在下面，我们令 $H$ 表示 $\e$ 上的一个幂幺代数群。（实践中 $H=\bN$。）

\pnumt{类与迹。}{假设 $\mathcal{C}\in\DGCat_{\mathrm{cont}}$ 是一个可对偶的 DG 范畴，$T\colon\mathcal{C}\to\mathcal{C}$ 是一个自函子。}

令 $\mathcal{C}^c\subseteq\mathcal{C}$ 表示紧对象子范畴，并令 $\mathcal{C}^{c,T,\mathrm{lax}}$ 表示由对 $(F,\alpha)$ 构成的范畴，其中 $F\in\mathcal{C}^c$ 且 $\alpha\colon F\to T(F)$ 是 $\mathcal{C}$ 中的态射。

给定如上某个 $(F,\alpha)$，存在一个典范点 $\cl(F,\alpha)\in\Om^\infty\tr_{\mathcal{C}}(T)$。事实上，这由 §3.2.5 中迹的函子性推出：给 $\Vect$ 配备恒等自映射，把 $F$ 看作函子 $\Vect\to\mathcal{C}$，把 $\alpha$ 看作 lax 缠绕映射，于是函子性给出映射 $\e=\tr_{\Vect}(\id_{\Vect})\to\tr_{\mathcal{C}}(T)\in\Vect$，即一个点 $\cl(F,\alpha)\in\Om^\infty\tr_{\mathcal{C}}(T)$。

更一般地，我们回忆存在一个映射 $\cl(-)\colon K(\mathcal{C}^{c,T,\mathrm{lax}})\to\tr_{\mathcal{C}}(T)\in\Spectra$，从 $\mathcal{C}^{c,T,\mathrm{lax}}$ 的 $K$-理论谱到 $T$ 的迹（后者经由遗忘函子 $\Vect\to\Spectra$ 视为谱）。

\begin{notation}
注意 $\Vect^{c,\id,\mathrm{lax}}=\{(W,\phi):W\in\Vect^c,\ \phi\colon W\to W\}$ 是一个对称幺半范畴，因此在上述设定中典范地作用在 $\mathcal{C}^{c,T,\mathrm{lax}}$ 上。显式地，对 $(W,\phi)\in\Vect^{c,\id,\mathrm{lax}}$ 与 $(F,\alpha)\in\mathcal{C}^{c,T,\mathrm{lax}}$，$W\otimes F$ 配备自同态 $\phi\otimes\id+\id\otimes\alpha$。

在类映射下，有：
\begin{equation}
  \cl(W\otimes F,\phi\otimes\id+\id\otimes\alpha)=\tr_W(\phi)\cdot\cl(F,\alpha).
\end{equation}
（我们并不需要这一点，但这个等式可以很容易地提升为谱层面上的适当陈述。）
\end{notation}

\pnumt{问题的陈述。}{回忆 $H$ 是幂幺的。由 [2] 命题 3.3.2，$\LS^{\restr}_{H}$ 是拟紧代数叠；特别地，其结构层 $\OO_{\LS^{\restr}_{H}}\in\QCoh(\LS^{\restr}_{H})$ 是紧的（与非幂幺群的情形不同）。按滥用记号，我们仍用 $\OO_{\LS^{\restr}_{H}}$ 表示在完全忠实嵌入
$\QCoh(\LS^{\restr}_{H})\hookrightarrow\IndCoh(\LS^{\restr}_{H})$ 下 $\IndCoh(\LS^{\restr}_{H})$ 中的“同一个”对象（在 $\IndCoh$ 文献中通常记作“$\Xi$”）。}

我们有一个映射 $\tau_0\colon\mathbf{B}H\to\LS^{\restr}_{H}$，对应于平凡局部系统。于是可以构造
\[
  \tau^{\IndCoh}_{0,*}(\OO_{\mathbf{B}H})\in\IndCoh(\LS^{\restr}_{H}).
\]
注意 $\OO_{\LS^{\restr}_{H}}$ 与 $\tau^{\IndCoh}_{0,*}(\OO_{\mathbf{B}H})$ 这两个对象都是凝聚的，并带有显然的典范 Frobenius 等变结构。因此，我们可以构造它们的类：\footnote{由 $H$ 的幂幺性，注意 $\LS^{\arthm}_{H}$ 的每个点都是非共振的，即映射 $H^{\ad}/H\to\LS^{\arthm}_{H}$ 是同构。}
\[
  \cl(\OO_{\LS^{\restr}_{H}},\alpha),\quad
  \cl(\tau^{\IndCoh}_{0,*}(\OO_{\mathbf{B}H}),\alpha)
  \in\Om^\infty\Gamma(\LS^{\arthm}_{H},\omega).
\]
我们的目标是比较这两个类。

\pnumt{一个 $L$-值。}{令 $\zeta_X(t)=\sum_{n\ge 0}|X_0(\mathbf{F}_{q^n})|t^n\in\Q(t)\subseteq\Q((t))$ 表示曲线 $X$ 的 $\zeta$ 函数。}

我们提醒 $\zeta$ 函数具有形式：
\[
  \zeta_X(t)=\frac{p_X(t)}{(1-t)(1-qt)},\qquad p_X(t)\in\Q[t].
\]
令 $\zeta^{\star}_X(t)=(1-t)\cdot\zeta_X(t)$。由 Weil 猜想，$\zeta^{\star}_X(1)$ 非零，因而等于 $\zeta_X(t)$ 在 $t=1$ 处 Taylor 展开\footnote{注意——与数论中不同——我们是在变量 $t=q^{-s}$ 中展开，而不是在 $s$ 本身中展开。}的首项。

\pnumt{主引理。}{我们将证明：}

\begin{lem}
存在一个等价
\[
  \cl(\tau^{\IndCoh}_{0,*}(\OO_{\mathbf{B}H}),\alpha)
  \simeq \zeta^{\star}_X(1)^{\dim H}\cdot\cl(\OO_{\LS^{\restr}_{H}},\alpha)
  \in\Om^\infty\Gamma(\LS^{\arthm}_{H},\omega).
\]
\end{lem}

撇开同伦的信息，这个结果只是说上述两个点在集合
$\pi_0(\Om^\infty\Gamma(\LS^{\arthm}_{H},\omega))=H^0\Gamma(\LS^{\arthm}_{H},\omega)$ 中的像相等。

\begin{proof}
\textbf{第 1 步。}我们从一个玩具模型开始。

令 $V$ 与 $W$ 是配备自同态 $\phi_V$ 与 $\phi_W$ 的有限维向量空间。令 $\mathcal{Y}$ 表示叠（在 $\e$ 上）$V\times\Om_0W$，我们把它配备自映射 $\phi=\phi_V\times\Om_0\phi_W$；这里 $\Om_0W$ 是导出环路空间 $0\times_W 0$。

例如我们令 $\OO_0\in\Coh(\mathcal{Y})$ 表示原点的结构层，$\OO_V\in\Coh(\mathcal{Y})$ 表示 $V\subseteq\mathcal{Y}$ 的结构层，等等。

我们使用 $\mathcal{Y}$ 上 lax $\phi_*$-等变凝聚层范畴 $\Coh(\mathcal{Y})^{\phi_*,\mathrm{lax}}$，即
$\Coh(\mathcal{Y})^{\phi_*,\mathrm{lax}}=\{F\in\Coh(\mathcal{Y}),\ \alpha\colon F\to\phi_*(F)\}$（见 §4.6.1）。

Koszul 分解给出恒等式：
\begin{equation}
\begin{gathered}
  \sum(-1)^i[\Lambda^iW^{\vee}\otimes\OO_V]\simeq[\OO_{\mathcal{Y}}],\\
  \sum(-1)^j[\Lambda^jV^{\vee}\otimes\OO_V]\simeq[\OO_0].
\end{gathered}
\end{equation}
这里记号的含义如下。首先，$\Lambda^iW^{\vee}\otimes\OO_{\mathcal{Y}}\in\Coh(\mathcal{Y})^{\phi_*,\mathrm{lax}}$ 配备来自记号 4.6.1.1 的 lax 等变结构，其中 $\Lambda^iW^{\vee}$ 配备自同态 $\Lambda^i\phi_W^{\vee}$；对 $\Lambda^jV^{\vee}\otimes\OO_V$ 有类似的记号。记号 $[-]$ 用于表示 $K$-理论谱\footnote{对我们的目的而言，也可以简单地使用 Grothendieck 群。}$K(\Coh(\mathcal{Y})^{\phi_*,\mathrm{lax}})$ 中的类。

由 (4.10)，我们得到：
\[
\begin{gathered}
  \cl(\OO_0,\alpha)=\sum(-1)^j\tr(\Lambda^j\phi_V^{\vee})\cdot\cl(\OO_V,\alpha)
  =\sum(-1)^j\tr(\Lambda^j\phi_V)\cdot\cl(\OO_V,\alpha)\\
  =\det(\id_V-\phi_V)\cdot\cl(\OO_V,\alpha)\in\Om^\infty\Gamma(\mathcal{Y}^{\phi},\omega),
\end{gathered}
\]
其中 $\mathcal{Y}^{\phi}$ 表示导出不动点。

类似地，我们有：
\[
  \cl(\OO_{\mathcal{Y}},\alpha)=\det(\id_W-\phi_W)\cdot\cl(\OO_V,\alpha)
  \in\Om^\infty\Gamma(\mathcal{Y}^{\phi},\omega).
\]
比较这两个恒等式，我们得到：
\[
  \det(\id_W-\phi_W)\cdot\cl(\OO_0,\alpha)
  =\det(\id_V-\phi_V)\cdot\cl(\OO_{\mathcal{Y}},\alpha).
\]
现在假设 $\det(\id_W-\phi_W)$ 非零，从而在域 $\e$ 中可逆。我们得到：
\[
  \cl(\OO_0,\alpha)=\frac{\det(\id_V-\phi_V)}{\det(\id_W-\phi_W)}\cdot\cl(\OO_{\mathcal{Y}},\alpha)
  \in\Om^\infty\Gamma(\mathcal{Y}^{\phi},\omega).
\]
取 $V=H^1_{\et}(X)$ 与 $W=H^2_{\et}(X)$，配备它们各自的 Frobenius 自同态，我们注意到：
\[
  \frac{\det(\id_V-\phi_V)}{\det(\id_W-\phi_W)}=\zeta^{\star}_X(1),
\]
这正是 Grothendieck 迹公式。

\textbf{第 2 步。}其次，假设我们处于以下更一般的设定中。

我们假设 $\mathcal{Z}$ 是一个 QCA 叠，配备自同态 $\phi=\phi_{\mathcal{Z}}$，并且配备一个与映射 $\phi$ 缠绕的拟光滑映射 $\mathcal{Z}\to\mathcal{Y}=V\times\Om_0W$。

我们令 $\mathcal{Z}_0$ 表示 $\mathcal{Z}$ 在 $0\in\mathcal{Y}$ 上的纤维。我们注意到 $\mathcal{Z}_0$ 是最终余连通的，因此 $\OO_{\mathcal{Z}_0}$ 是 $\mathcal{Z}$ 上的凝聚层。

前面的分析于是表明：
\begin{equation}
  \cl(\OO_{\mathcal{Z}_0},\alpha)
  =\frac{\det(\id_V-\phi_V)}{\det(\id_W-\phi_W)}\cdot\cl(\OO_{\mathcal{Z}},\alpha)
  \in\Om^\infty\Gamma(\mathcal{Z}^{\phi},\omega)
\end{equation}
（假设 $1$ 不是 $\phi_W$ 的特征值）。

\textbf{第 3 步。}现在我们希望把上述形式体系应用于推导我们的断言。

选取一列嵌套子群 $\{1\}=H_0\subseteq H_1\subseteq\cdots\subseteq H_r=H$，其中每个 $H_i$ 在 $H$ 中正规，且 $H_{i+1}/H_i\simeq\mathbb{G}_a$。我们指出 $r=\dim(H)$。

然后定义代数叠：
\[
  \mathcal{Z}_i:=\LS^{\restr}_{H}\times_{\LS^{\restr}_{H/H_i}}\mathbf{B}(H/H_i),
\]
其中 $\mathbf{B}(H/H_i)\to\LS^{\restr}_{H/H_i}$ 是映射 $\tau_0$，即它对应于 $X$ 上的平凡 $H/H_i$-局部系统。注意 $\mathcal{Z}_r=\LS^{\restr}_{H}$，$\mathcal{Z}_0=\mathbf{B}H$，并且我们有闭嵌入：
\[
  \mathcal{Z}_0\hookrightarrow\mathcal{Z}_1\hookrightarrow\cdots\hookrightarrow\mathcal{Z}_r.
\]
我们可以把 $\mathcal{Z}_i$ 的定义改写如下。

注意 $H/H_i$ 作用在分类叠 $\mathbf{B}H_i$ 上；形式地说，这由纤维序列 $\mathbf{B}H_i\to\mathbf{B}H\to\mathbf{B}(H/H_i)$ 编码。展开定义，这诱导了 $H/H_i$ 在 $\LS^{\restr}_{H_i}$ 上的一个作用。于是我们有：
\[
  \mathcal{Z}_i\simeq\LS^{\restr}_{H_i}/(H/H_i).
\]
现在对每个 $i$，我们有一个图：
\[
\begin{tikzcd}[column sep=1.6em, row sep=1.4em]
  H/H_i \arrow[r] \arrow[d, Rightarrow]
    & H/H_{i+1} \arrow[r] \arrow[d, Rightarrow]
    & \mathbf{B}G_a \arrow[d, Rightarrow]\\
  \LS^{\restr}_{H_i} \arrow[r]
    & \LS^{\restr}_{H_{i+1}} \arrow[r]
    & \LS^{\restr}_{G_a}
\end{tikzcd}
\]
其中各行是纤维序列，且顶行是群的纤维序列。这里 $\mathbf{B}G_a$ 在 $\LS^{\restr}_{G_a}$ 上的作用由群叠的同态 $\mathbf{B}G_a\to\LS^{\restr}_{G_a}$ 诱导，后者对应于沿 $X\to\Spec(k)$ 拉回局部系统（即 $G_a$ 的映射 $\tau_0$）。在这个图中取商，并利用\footnote{我们指出，由纯度，存在一个与 Frobenius 相容的典范这样的分裂。}
\[
  \LS^{\restr}_{G_a}\simeq\mathbf{B}G_a\times H^1_{\et}(X)\times\Om_0H^2_{\et}(X),
\]
我们得到一个纤维方块：
\[
\begin{tikzcd}[column sep=2.2em, row sep=1.6em]
  \mathcal{Z}_i \arrow[r] \arrow[d] & \mathcal{Z}_{i+1} \arrow[d] & \\
  \Spec(k) \arrow[r] & \LS^{\restr}_{G_a}/\mathbf{B}G_a \arrow[r, equal]
    & H^1_{\et}(X)\times\Om_0H^2_{\et}(X)
\end{tikzcd}
\]
现在由前一步的归纳，我们得到结果。\footnote{形式地说，归纳应当在 $K$-理论类上进行，并推广 (4.11)。我们只在最后才映射到 Hochschild 同调。}
\end{proof}

\pnumt{一段扩展的题外话：经由范畴迹处理发散级数。}{我们解释思考上述引理 4.6.4.1 证明的一个一般格式。这些材料是非正式的，可以跳过。不过，我们相信这是一个重要的哲学要点，值得强调。}

粗略地说，想法是：所谓的\emph{范畴泛函分析}（例如关于 $\Perf$ 与 $\Coh$ 之间区别的精细考虑）通过范畴迹与实际分析（例如对无穷级数求和）相关联。引人注目的是，我们将看到 Hochschild 同调使我们有时能够“正确地”计算无穷和，而完全不必提及这些和所在的域 $\e$ 上的任何拓扑。

我们考虑以下几何设定。令 $\mathcal{Y}$ 是（$\e$ 上的）代数叠，我们假设它拟光滑且 QCA。假设 $\mathcal{Y}$ 配备自映射 $\phi\colon\mathcal{Y}\to\mathcal{Y}$。函子
\[
\begin{gathered}
  \Upsilon_{\mathcal{Y}}\colon\QCoh(\mathcal{Y})\to\IndCoh(\mathcal{Y}),\\
  F\mapsto F\otimes_{\OO_{\mathcal{Y}}}\omega_{\mathcal{Y}}
\end{gathered}
\]
保持紧对象，并缠绕源与靶的自映射 $\phi^*$ 与 $\phi^!$。此外，这个函子是 $\QCoh(\mathcal{Y})$-模范畴的态射。

现在回忆（例如 [15] §3.8.8），对于一个可对偶的 $\QCoh(\mathcal{Y})$-模范畴 $\mathcal{C}$ 与一个与 $\phi^*$ 适当相容的自函子 $T\colon\mathcal{C}\to\mathcal{C}$，存在一个典范对象
\[
  \tr^{\mathrm{enh}}(T)\in\QCoh(\mathcal{Y}^{\phi}),
\]
其基本性质是 $\Gamma(\mathcal{Y},\tr^{\mathrm{enh}}(T))=\tr(T)\in\Vect$。这一构造满足迹的通常函子性。我们有\footnote{在这个公式中，我们把 $\omega_{\mathcal{Y}^{\phi}}$ 看作 $\QCoh$ 的对象，而不是 $\IndCoh$ 的对象。换言之，我们隐含地取 $\IndCoh$ 中的“真正”对偶化层，并对它应用遗忘函子 $\Psi\colon\IndCoh(\mathcal{Y}^{\phi})\to\QCoh(\mathcal{Y}^{\phi})$。由于 $\omega_{\mathcal{Y}^{\phi}}$ 是线丛，这是一件相当温和的事情，所以我们不专门在记号中标明它。}如下公式：
\[
  \tr^{\mathrm{enh}}(\phi^*)=\OO_{\mathcal{Y}^{\phi}},\qquad
  \tr^{\mathrm{enh}}(\phi^!)=\omega_{\mathcal{Y}^{\phi}},
\]
于是 $\Upsilon_{\mathcal{Y}}$ 给出一个典范映射
\[
  \tau\colon\OO_{\mathcal{Y}^{\phi}}=\tr^{\mathrm{enh}}(\phi^*)\to\tr^{\mathrm{enh}}(\phi^!)=\omega_{\mathcal{Y}}
  \in\QCoh(\mathcal{Y}^{\phi}).
\]
（这个映射 $\tau$ 可以看作导出不动点上的一个弱 Calabi--Yau 结构。）

我们令 $\mathcal{Y}^{\phi,\mathrm{good}}\subseteq\mathcal{Y}^{\phi}$ 表示 $\tau$ 为同构的点的轨迹。注意 $\mathcal{Y}^{\phi,\mathrm{good}}$ 包含 $(\mathcal{Y}^{\mathrm{sm}})^{\phi}$（$\mathcal{Y}$ 的光滑轨迹的不动点），但一般更大：事实上可以验证 $\mathcal{Y}^{\phi,\mathrm{good}}$ 恰是 $\mathcal{Y}^{\phi}$ 的拟光滑轨迹。

现在，任何带自映射 $\alpha\colon F\to\phi^*(F)$ 的完美对象 $F\in\Perf(\mathcal{Y})$ 给出一个类
\[
  \cl(F,\alpha)_{\QCoh(\mathcal{Y})}\in\Om^\infty\Gamma(\mathcal{Y}^{\phi},\OO),
\]
即 $\phi$ 的不动点 $\mathcal{Y}^{\phi}$ 上的一个函数。在这个记号中，我们用下标 $\cl(-,-)_{\QCoh}$ 强调我们把 $F$ 看作 $\QCoh$ 的对象（这很快就会成为一个重要的区分）。这个函数可以相当显式地理解：在点 $y\in\mathcal{Y}^{\phi}$ 处，我们取所得映射的迹：
\begin{equation}
  \beta_y\colon y^*(F)\xrightarrow{\alpha}y^*(\phi^*(F))=\phi(y)^*(F)\cong y^*(F),
\end{equation}
其中最后一个同构使用了 $y$ 是不动点所隐含的等同 $y\simeq\phi(y)$。换言之，我们有：
\[
  \cl(F,\alpha)_{\QCoh(\mathcal{Y})}=(y\mapsto\tr(\beta_y)).
\]
现在改为假设 $F\in\Coh(\mathcal{Y})$，但仍配备映射 $\alpha\colon F\to\phi^*(F)$。因为 $F$ 在 $\QCoh(\mathcal{Y})$ 中可能不是紧的，我们不能再构造它在 $\Gamma(\mathcal{Y}^{\phi},\OO)$ 中的类。不过，我们可以扭一下并构造 $F\otimes\omega_{\mathcal{Y}}$，它属于 $\Coh$，因为 $\mathcal{Y}$ 拟光滑（从而 Gorenstein）。于是我们得到映射
\[
  \widetilde\alpha=\alpha\otimes\id\colon
  F\otimes\omega_{\mathcal{Y}}\to\phi^*(F)\otimes\omega_{\mathcal{Y}}
  =\phi^!(F\otimes\omega_{\mathcal{Y}}).
\]
因此，我们可以构造类
\[
  \cl(F\otimes\omega_{\mathcal{Y}},\widetilde\alpha)_{\IndCoh}
  \in\Gamma(\mathcal{Y}^{\phi},\omega_{\phi}).
\]
显然地，在 $F\in\Perf(\mathcal{Y})\subseteq\Coh(\mathcal{Y})$ 的特殊情形，我们有：
\begin{equation}
  \cl(F\otimes\omega_{\mathcal{Y}},\widetilde\alpha)_{\IndCoh}
  =\tau(\cl(F,\alpha)_{\QCoh}).
\end{equation}
遵循这个等式，我们定义\emph{正则化类}
\[
  \cl^{\mathrm{reg}}(F,\alpha)_{\QCoh}\in\Om^\infty\Gamma(\mathcal{Y}^{\phi,\mathrm{good}},\OO),
\]
为 $\cl(F\otimes\omega_{\mathcal{Y}},\widetilde\alpha)_{\IndCoh}$ 在如下复合下的像：
\[
  \Gamma(\mathcal{Y}^{\phi},\omega)\to
  \Gamma(\mathcal{Y}^{\phi,\mathrm{good}},\omega)
  \xrightarrow{\tau^{-1}}\Gamma(\mathcal{Y}^{\phi,\mathrm{good}},\OO).
\]
由构造，当 $F$ 完美时，这个正则化类与（$\cl(F,\alpha)_{\QCoh}$ 到 $\mathcal{Y}^{\phi,\mathrm{good}}$ 的限制）重合。

假设 $y\in\mathcal{Y}^{\phi}$。映射 (4.13) 仍有意义。然而，如果 $F$ 在 $y$ 附近不完美，那么虽然复形 (4.13) 在每个次数都是有限维的，它却下方无界，因此 $\beta_y$ 的迹没有良好定义。我们定义\emph{正则化迹}
\[
  \tr^{\mathrm{reg}}(\beta_y):=(\cl^{\mathrm{reg}}(F,\alpha)_{\QCoh})(y)
\]
为正则化类在 $y\in\mathcal{Y}^{\phi}$ 处的值。

启发式地，正则化迹可以看作给无穷和赋予一个实际值：
\begin{equation}
  \sum_{i\in\Z}(-1)^i\tr(H^i(\beta_y))\colon H^i(y^*(F))\to H^i(y^*(F))\in\e,
\end{equation}
其中我们重申每个求和项都有良好定义，求和项对 $i\gg 0$ 消失，但一般而言会出现无穷多项。

\begin{example}
我们解释在最简单的情形中这是如何运作的。假设 $\mathcal{Y}=\Om_0\A^1$，$\phi$ 是乘以数 $\lambda\in\e$。取层 $F$ 为 $\OO_0$，即点 $0\in\Om_0\A^1$ 的结构层。注意 $0$ 典范地是 $\phi$ 的不动点，所以我们可以把 $0$ 看作 $\mathcal{Y}^{\phi}$ 的点。由标准计算，$0^*(\OO_0)$ 在偶的非正上同调次数有 $1$ 维上同调，在这些次数之外上同调消失；此外，映射 $\beta_0$ 在 $H^{-2i}(0^*(\OO_0))$ 上作为乘以 $\lambda^i$ 作用。因此，(4.15) 中的和是几何级数 $\sum_{i\ge 0}\lambda^i$。我们强调这是一个形式表达式；目前，$\lambda$ 是域 $\lambda$ 的任意元素，绝不是“小”的。% [Tutor Note: 原文此处作 the field λ，应为系数域 e；疑为笔误，照录并登记。]

现在假设 $\lambda\neq 1$。则 $0\in\mathcal{Y}^{\phi,\mathrm{good}}$（事实上 $\mathcal{Y}^{\phi,\mathrm{good}}=\mathcal{Y}^{\phi}=\Spec(\e)=\{0\}$）。于是正则化迹 $\tr^{\mathrm{reg}}(\beta_0)$ 有良好定义，而引理 4.6.4.1 证明中第 1 步的（完全初等的）论证算出
\[
  \tr^{\mathrm{reg}}(\beta_0)=\frac{1}{1-\lambda}.
\]
换言之，我们给出了几何级数公式 $\sum\lambda^i\text{ “=” }\frac{1}{1-\lambda}$ 的直接的、纯粹代数的意义，而通常这需要我们知道 $\lambda^i$ 在某种适当的解析意义下 $i\to\infty$ 时趋于 $0$。
\end{example}

\begin{remark}[正则化迹与函数方程]
现在假设 $\phi\colon\mathcal{Y}\to\mathcal{Y}$ 事实上是同构。则 $\phi^*=\phi^!$（比如，作为限制在 $\Perf$ 或 $\Coh$ 上的函子）。因此，对于带 $\alpha\colon F\to\phi^*(F)$ 的 $F\in\Perf(\mathcal{Y})$，我们也得到映射 $\alpha'\colon F\to\phi^!(F)$。对 $y\in\mathcal{Y}^{\phi}$，我们得到典范映射
\[
  \gamma_y\colon y^!(F)\to y^!(F),
\]
以与 $\beta_y$ 相同的方式定义。由于 $F$ 完美，我们有
\[
  y^!(F)=y^!(F\otimes\omega_{\mathcal{Y}}\otimes\omega_{\mathcal{Y}}^{-1})
  =y^!(F\otimes\omega_{\mathcal{Y}})\otimes y^*(\omega_{\mathcal{Y}}^{-1})
  =y^*(F)\otimes y^*(\omega_{\mathcal{Y}})^{\otimes-1}.
\]
这个映射把 $\gamma_y$（对 $F$）与 $\beta_y$（对 $F$ 与 $\omega_{\mathcal{Y}}$ 两者）缠绕起来。如果我们令 $\epsilon_y$ 为映射
\[
  y^*(\omega_{\mathcal{Y}})\to y^*(\omega_{\mathcal{Y}})
\]
的迹（用 $\gamma_y$ 与显然的同构 $\omega_{\mathcal{Y}}\simeq\phi^!(\omega_{\mathcal{Y}})=\phi^*(\omega_{\mathcal{Y}})$ 构造），则我们发现
\[
  \epsilon_y\cdot\tr(\gamma_y)=\tr(\beta_y).
\]
现在我们可以完全像对凝聚 $F$ 那样，当 $y\in\mathcal{Y}^{\phi,\mathrm{good}}$ 时定义 $\tr^{\mathrm{reg}}(\gamma_y)$。我们得到一个同义反复的“函数方程”：
\[
  \epsilon_y\cdot\tr^{\mathrm{reg}}(\gamma_y)=\tr^{\mathrm{reg}}(\beta_y)
\]
（其中 $\epsilon_y$ 被看作一个 $\epsilon$-因子）。

让我们看看这个逻辑在例 4.6.5.1 的设定中如何展开。我们应当有 $\lambda\neq 0$，使得 $\phi$ 是同构。我们注意到 $0^!(\OO_0)$ 在偶的非负次数有上同调，而 $\gamma_0$ 在 $H^{2i}(0^!(\OO_0))$ 上的作用是乘以 $\lambda^{-i}$。此外，$0^*(\omega)=\e[-1]$，其中“$\gamma$”算子是乘以 $\lambda^{-1}$。因此，正则化迹 $\tr^{\mathrm{reg}}(\gamma_0)$ 启发式地使和
\[
  \tr^{\mathrm{reg}}(\gamma_0)\text{ “=” }\sum_{i\ge 0}\lambda^{-i}
\]
有意义，而 $\epsilon_0=-\lambda^{-1}$。

因此，在这个情形，我们的函数方程启发式地给出：
\[
  -\lambda^{-1}\sum_{i\ge 0}\lambda^{-i}\text{ “=” }
  \epsilon_0\tr^{\mathrm{reg}}(\gamma_0)=\tr^{\mathrm{reg}}(\beta_0)
  \text{ “=” }\sum_{i\ge 0}\lambda^i.
\]
\end{remark}

这个所得方程 $-\lambda^{-1}\sum_{i\ge 0}\lambda^{-i}\text{ “=” }\sum_{i\ge 0}\lambda^i$ 是发散级数研究中的一个经典式子；在 $\mathbb{C}$ 上，左边对 $|\lambda|>1$ 有定义，而右边对 $|\lambda|<1$ 有定义，但当然，这两个函数的解析延拓在其定义域上重合。

我们注意到，这类对发散级数的操作与代数曲线 $\zeta$ 函数的函数方程密切相关。

\begin{remark}
用上述思想的某种版本，在范畴上解释函数域上自守形式解析理论的某些解析方面，将是非常有意义的。
\end{remark}

\pnumt{变体。}{在实践中，我们需要 §4.6.4 讨论的一个轻微推广。}

像之前一样令 $H$ 是一个幂幺群，现在令 $S$ 是一个通过自同构作用在 $H$ 上的环面。（实践中，$S=\bT$ 作用在 $H=\bN$ 上。）为简洁起见，我们令 $Q$ 表示半直积 $S\ltimes H$。

在这个情形，我们定义有理映射
\[
  \zeta^{\star}_{X,H,S}\colon S\dashrightarrow\A^1
\]
通过公式
\[
  (s\in S)\mapsto
  \frac{\det\big(\id-\Ad_{s^{-1}}\otimes\phi_X\curvearrowright\mathfrak{h}\otimes H^1_{\et}(X)\big)}
       {\det\big(\id-\Ad_{s^{-1}}\otimes\phi_X\curvearrowright\mathfrak{h}\otimes H^2_{\et}(X)\big)}.
\]
这里 $\phi_X$ 是作用在 $X$ 的 \et{} 上同调上的 Frobenius，而我们在让 $\Ad_{-}$ 表示由 $S$ 在 $H$ 上的作用导出的 $S$ 在 $\mathfrak{h}$ 上的作用时，稍微滥用了记号。

\begin{remark}
注意 $\zeta^{\star}_{X,H,S}$ 在 $1\in S$ 处有定义，并在该处取值 $\zeta^{\star}_X(1)^{\dim H}$（见引理 4.6.4.1）。
\end{remark}

\begin{remark}
假设 $\mu_1,\dots,\mu_r\colon S\to\mathbb{G}_m$ 是 $S$ 在其表示 $\mathfrak{h}$ 中出现的特征标，按重数计算（故 $r=\dim(H)$）。则我们有：
\[
  \zeta^{\star}_{X,H,S}(s)=\prod_{i=1}^{r}\zeta^{\star}_X(\mu_i(s^{-1})).
\]
特别地，$\zeta^{\star}_{X,H,S}$ 的定义域是 $\bigcap_i\{s\in S\mid\mu_i(s)\neq q\}$，而 $\zeta^{\star}_{X,H,S}$ 在 $\bigcap_i\{s\in S\mid\mu_i(s)\notin\mathcal{R}_X\}$ 上（有定义且）可逆。
\end{remark}

同时引入以下记号会很方便：
\[
\begin{gathered}
  p_{X,H,S}(s)=\det\big(\id-\Ad_{s^{-1}}\otimes\phi_X\curvearrowright\mathfrak{h}\otimes H^1_{\et}(X)\big)
  =\prod_{i=1}^{r}p_X(\mu_i(s^{-1})),\\
  q_{X,H,S}(s)=\det\big(\id-\Ad_{s^{-1}}\otimes\phi_X\curvearrowright\mathfrak{h}\otimes H^2_{\et}(X)\big)
  =\prod_{i=1}^{r}(1-q\cdot\mu_i(s^{-1})),
\end{gathered}
\]
于是 $p_{X,H,S}$ 与 $q_{X,H,S}$ 是 $S$ 上的（正则）函数，且
$\zeta^{\star}_{X,H,S}=\dfrac{p_{X,H,S}}{q_{X,H,S}}$。

我们引入记号：
\[
  \LS^{\restr}_{Q;S}:=\LS^{\restr}_{Q}\times_{\LS^{\restr}_{S}}\mathbf{B}S
  =(\LS^{\restr}_{H})/S.
\]
我们观察到 $\LS^{\restr}_{Q;S}$ 是拟紧代数叠（由 $H$ 的幂幺性）。我们令 $\LS^{\arthm}_{Q;S}$ 表示 $\LS^{\restr}_{Q;S}$ 的 Frobenius 不动点。显式地，我们有：
\[
  \LS^{\arthm}_{Q;S}=\LS^{\arthm}_{Q}\times_{\LS^{\arthm}_{S}}S^{\ad}/S.
\]
我们令 $\tau_1\colon\mathbf{B}Q\to\LS^{\restr}_{Q;S}$ 表示显然的映射。

\begin{remark}
更显式地说，我们提醒 $S^{\ad}/S\xrightarrow{\sim}\LS^{\arthm,\circ}_{S}$（由定理 4.3.3.1），并回忆 $S$ 是环面。因此，$\LS^{\arthm}_{Q;S}$ 是 $\LS^{\arthm}_{Q}$ 中单位元的连通分支。
\end{remark}

\begin{lem}
存在一个等价
\[
  q_{X,H,S}\cdot\cl(\tau^{\IndCoh}_{1,*}(\OO_{\mathbf{B}Q}),\alpha)
  \simeq p_{X,H,S}\cdot\cl(\OO_{\LS^{\restr}_{Q;S}},\alpha)
  \in\Om^\infty\Gamma(\LS^{\arthm}_{Q;S},\omega).
\]
\end{lem}

\begin{remark}
非正式地，这个引理应被理解为说：
\[
  \cl(\tau^{\IndCoh}_{1,*}(\OO_{\mathbf{B}Q}),\alpha)
  \simeq\zeta^{\star}_{X,H,S}\cdot\cl(\OO_{\LS^{\restr}_{Q;S}},\alpha).
\]
\end{remark}

\begin{proof}[引理 4.6.6.4 的证明]
其证明本质上与引理 4.6.4.1 的证明相同。差别如下。

第一，在第 1 步中，应当假设 $V$ 与 $W$ 是 $S$-表示，并应当在 (4.11) 中计入 $S$-作用。对偶表示 $V^{\vee}$ 与 $W^{\vee}$ 出现在 (4.11) 中，这解释了上文 $\zeta^{\star}_{X,H,S}$ 的定义中出现 $s^{-1}$ 而非 $s$。

第二，应当注意引理 4.6.4.1 证明中的子群 $H_i$ 可以取为在 $S$-作用下不变的（证明：把 $S$ 在 $\mathfrak{h}/[\mathfrak{h},\mathfrak{h}]$ 上的作用对角化，然后归纳）。

除此之外，论证逐字相同。
\end{proof}
